\documentclass[11pt,a4paper]{article}

\usepackage{mathptmx}
\usepackage[utf8]{inputenc}
\usepackage[margin=1.0in]{geometry}
\usepackage{amsmath}
\usepackage{amssymb}
\usepackage{amsthm}
\usepackage{booktabs}
\usepackage{array}
\usepackage[hidelinks]{hyperref}

% --- compact but standard layout ---
\setlength{\emergencystretch}{3em}
\hyphenpenalty=200
\setlength{\tabcolsep}{4.5pt}
\renewcommand{\arraystretch}{0.97}
\let\oldquote\quote
\let\endoldquote\endquote
\renewenvironment{quote}{\begin{oldquote}\small}{\end{oldquote}}
\makeatletter
\renewcommand\section{\@startsection{section}{1}{\z@}%
  {-2.2ex \@plus -0.6ex \@minus -.2ex}{1.4ex \@plus.2ex}{\normalfont\large\bfseries}}
\renewcommand\subsection{\@startsection{subsection}{2}{\z@}%
  {-2.0ex\@plus -0.6ex \@minus -.2ex}{1.0ex \@plus .2ex}{\normalfont\normalsize\bfseries}}
\makeatother

\newtheorem{statement}{Statement}
\newtheorem{definition}[statement]{Definition}
\newtheorem{proposition}[statement]{Proposition}
\newtheorem{lemma}[statement]{Lemma}
\newtheorem{theorem}[statement]{Theorem}
\newtheorem{corollary}[statement]{Corollary}

\newcommand{\MAX}{\mathrm{MAX}}
\newcommand{\MAXTWO}{\mathrm{MAX2}}
\newcommand{\nc}{\mathrm{nc}}
\newcommand{\LCA}{\mathrm{LCA}}
\newcommand{\lb}{\mathrm{lp}}

\title{Twenty tabulated lower bounds for sorting networks rest on
an invalid proof: an audit of Van Voorhis (1972), and a partial repair}

\author{}
\date{}

\begin{document}
\maketitle
\vspace{-3.5em}

\begin{abstract}
\noindent
The lower bounds on sorting-network size tabulated for $N \ge 13$ in the most
consulted public table rest on a two-value pruning theorem of Van Voorhis (1972)
whose published proof is invalid. Its
equations (5) and (6) are false: we exhibit a size-optimal three-comparator
network on which (5) overstates the quantity it claims to compute and (6)'s
Kraft sum equals $5/4$, and equation (7), the only route from them to the
result, rests on them alone. The false premise is that the MAX subnetwork of an
$N$-sorter consists of $N-1$ comparators, which ignores \emph{pass-through}
comparators, present in $61.0\,\%$ of an audited population. The same chain
reproduces the published table at every $N = 13,\dots,32$, so twenty entries
depend on it, and the two most consulted public tables contradict each other at
$N = 13,\dots,20$. We prove the theorem for two classes of networks, covering
$22.4\,\%$ of an audited corpus; half of that corpus has no proof by any route
given here. The best published bound, $S(13) \ge 43$, is untouched.
\end{abstract}

\section{Introduction}

A \emph{comparator} $(a,b)$ with $a<b$ on $n$ channels replaces the values on
channels $a$ and $b$ by their minimum on $a$ and their maximum on $b$. An
\emph{$n$-sorter} is a sequence of comparators that sorts every input, and
$S(n)$ is the least size of one --- the Bose--Nelson sorting problem
\cite{BN62}. $S(n)$ is known exactly only for $n \le 12$:

\begin{center}
\begin{tabular}{lccccccccccccc}
\toprule
$n$ & 1 & 2 & 3 & 4 & 5 & 6 & 7 & 8 & 9 & 10 & 11 & 12 & 13 \\
$S(n)$ & 0 & 1 & 3 & 5 & 9 & 12 & 16 & 19 & 25 & 29 & 35 & 39 & \textbf{?} \\
\bottomrule
\end{tabular}
\end{center}

\noindent
Values through $n=8$ are due to Floyd and Knuth \cite{FK73}, $S(9)$ and $S(10)$
to Codish, Cruz-Filipe, Frank and Schneider-Kamp \cite{CCFS16}, and $S(11)$ and
$S(12)$ to Harder \cite{Ha20}, each of the last four by exhaustive search with a
machine-checked or independently checked certificate \cite{CLS17,CS15,HaCert}.
At $n=13$ the best construction remains Juill\'e's 45-comparator network
\cite{Ju95}.

Two lower bounds for $S(13)$ are in circulation, and they differ. The
\textbf{one-value bound} $S(n) \ge S(n-1) + \lceil \log_2 n\rceil$ gives
$S(13) \ge 43$. It is Van Voorhis's \emph{IEEE Transactions} note \cite{VV72a},
Knuth's exercise 5.3.4--42 with a three-line answer \cite{Kn73}, and Harder's
Lemma 17, proved there in full \cite{Ha20}. This is the best \emph{published}
lower bound, and nothing here disturbs it (Section~\ref{sec:undisturbed}). The
\textbf{two-value bound} $S(N) \ge S(N-2) + P(2,N)$ with
$P(2,N) \ge \lceil \log_2 F(N)\rceil$ for an explicitly computable $F$ is the
content of Van Voorhis's chapter in \emph{Complexity of Computer Computations}
\cite{VV72}; with $F(13) = 392$ and the modern value $S(11) = 35$ it yields
$44$. The chapter's own Table~1 printed $L(13) = 42$, because in 1972 the best
available $S(11)$ was $33$, so the $44$ is a modern recombination. It entered
the public record on 2025-04-21 through a changelog line in Dobbelaere's table
\cite{Do25}; Wikipedia's table still prints $43$ \cite{Wik26}.

\medskip
\noindent
We audit the chapter's proof of the two-value bound and find it invalid. Two of
its structural steps are false, and the step that carries the whole conclusion is
derived from those two and from nothing else. The error is a single premise,
stated once on p.~121 and never revisited: that the MAX subnetwork of an
$N$-sorter consists of exactly $N-1$ comparators. It does not, whenever the
network contains what we call a \emph{pass-through} comparator, and
pass-throughs are ordinary: $61.0\,\%$ of an audited population of sorters has
one. The counterexample is a three-comparator network on three channels, and can
be checked by hand. We write this out because the argument
is fifty-four years old, has never been re-derived, re-proved or
machine-checked, and now supports a published table that people use.

Because the same inequality chain is what produces the published table above
$N=12$, the audit does not bear on $S(13)$ alone. We reconstruct the chain to
$N=32$ and it reproduces the published lower bounds at every one of
$N = 13,\dots,32$, exactly, twenty entries out of twenty --- the first written
derivation of a column published with no proof and no citation. Every one of
those twenty strictly exceeds the one-value floor, so every one depends on the
equation the audit breaks. The public record is meanwhile three-way
inconsistent: Wikipedia prints the one-value column, Dobbelaere prints the
two-value column, and OEIS A003075 \cite{OEIS}, credited in Dobbelaere's column
header, lists no lower bound at all for $N \ge 13$. This work says which of the
two columns is currently defensible.

We then repair what can be repaired. The missing step is an antichain argument;
we supply it for two classes of networks, prove a complementary bound at the
opposite extreme, and close one family of would-be repairs structurally. The
proved classes cover $22.4\,\%$ of the audited corpus, and the uncovered zone is
intermediate --- the inequality is provable at the pass-through-poor end and
again at the pass-through-rich end, with the open problem strictly in between.
The repair yields no new numerical bound at any $n$.

\paragraph{Related work.}
\emph{Exact values.} The four largest known values of $S(n)$ come from
exhaustive search with symmetry breaking \cite{CCFS16,Ha20}, formalized in Coq
for $n=9$ \cite{CLS17,CS15} and in Isabelle/HOL for $n = 11,12$ \cite{Ha20};
two-layer prefix generation \cite{CCS14} and depth-oriented pruning
\cite{BZ14,EM14,MG15} address parallel depth, which does not constrain size.
\emph{Lower bounds.} Apart from the information-theoretic
$\lceil \log_2 n! \rceil$ and the asymptotic $\Omega(n \log n)$ of
\cite{KLM95}, every published size lower bound above $n = 12$ is a chain from an
exact value by one of the two Van Voorhis inequalities.
\emph{Engagement with the chapter.} OpenAlex records 18 citing works for
\cite{VV72} as of 2026-08-22, two of them since 2024 \cite{Vlad25,CS24}; none
engages with equations (5)--(9), and the modern ones --- including the
reinforcement-learning constructions of \cite{Vlad25,BR26} at exactly these
widths --- cite the chapter for the one-value bound. A recent lower bound for \emph{stable standard} networks
\cite{TS26} concerns a different quantity. We found no erratum, correction, gap
note, reproof or formalization of the two-value argument, and no
correction-genre paper of any kind in the sorting-network literature.

\paragraph{Artifact, and what follows.}
Five standard-library Python programs accompany this paper; they are
deterministic, read-only and dependency-free, and re-run every number below in
under five minutes in total. The two that carry the audit were
written independently: one alongside the first reading of the chapter, the other
afterwards from the page images by a reader whose brief was to defend the
chapter. Every numerical claim below names the script that recomputes it;
check-level indexes, full proofs, refuted routes and known script defects are in
the artifact repository \cite{Art}.
Section~\ref{sec:prelim} fixes vocabulary; Section~\ref{sec:audit} is the audit;
Section~\ref{sec:record} is the literature record and the twenty entries;
Section~\ref{sec:repair} is the partial repair; Section~\ref{sec:undisturbed}
says what is not disturbed, including the unconditional structure of a
hypothetical minimum-size 13-sorter; Section~\ref{sec:trust} states what we
trust.

\section{Preliminaries}
\label{sec:prelim}

We match the verifier code, because several distinctions below are exactly what
the 1972 argument elides. Fix an $N$-sorter $T$; outputs ascend, so $o_N$ is the
last channel. A \emph{lead} is a wire segment: a channel's input segment, or a
comparator's output segment.

\begin{definition}[max-paths, branch tree, pass-throughs]
\label{def:one}
Give input channel $k$ a value strictly greater than all others. It leaves
every comparator it enters on the high output, tracing a unique \emph{max-path}
$p_k$ to $o_N$; put $\MAX(T) = \bigcup_k p_k$. Call a lead \emph{red} if some
one-hot max scenario puts the maximum on it, and \emph{blue} otherwise. A
traversed comparator is a \emph{branch node} if both its input leads are red,
and a \emph{pass-through} if exactly one is. The branch nodes form the
\emph{branch tree} $B$: $N$ leaves, $N-1$ internal nodes, rooted at the last
branch node before $o_N$. $T$ is \emph{clean} if it has no pass-through.
\end{definition}

For a branch node $c$ with branch-tree subtrees of heights $\lb(L(c))$ and
$\lb(R(c))$ and depth $\mathrm{depth}_B(c)$, put
\[
  a(c) := \lb(L(c)) + \lb(R(c)) + \mathrm{depth}_B(c) + 1,
  \qquad
  f(B) := \sum_{c} 2^{a(c)} .
\]
Then $f$ satisfies the chapter's Theorem~1 recursion, and $F(N) = \min_B f(B)$
reproduces its Table~1 exactly for $N \le 16$ under two independent
implementations (\texttt{verify\_huffman2.py},
\texttt{verify\_kraft\_dispute.py}).

For channels $x \ne y$, run the scenario \emph{maximum at $x$, second maximum at
$y$} and let $W(x,y)$ be the set of comparators either value traverses. Removing
them and rewiring leaves an $(N-2)$-sorter, so with
$p(2,T) := \max_{x \ne y} |W(x,y)|$ we have $|T| \ge S(N-2) + p(2,T)$: the
chapter's equation (2), not in dispute. Write
$W^*(c) = \max\{|W(x,y)| : \LCA_B(x,y) = c\}$, so
$p(2,T) = \max_c W^*(c)$. The disputed claim, the chapter's equation (8), is
\[
  p(2,T) \;\ge\; \lceil \log_2 f(B) \rceil . \tag{8}
\]

\begin{definition}[escapes; the hypotheses (C) and (D)]
\label{def:two}
In the scenario $(x,y)$ the two values first meet at $c = \LCA_B(x,y)$. $T$ is
\emph{escape-free} if the second maximum then occupies no red lead at or after
$c$'s low output; equivalently \cite{Art}, if it never re-meets the maximum at a
branch node. Write $nq(c)$ for the number of comparators the second maximum
traverses strictly after $c$, $ov(c)$ for those of them the maximum also
traverses, and $\gamma(c)$ for the number of pass-throughs on the maximum's stem
from $c$ to $o_N$. The two hypotheses isolated below are
\[
  \textbf{(C)}\quad \sum_c 2^{-nq(c)} \le 1,
  \qquad\qquad
  \textbf{(D)}\quad \gamma(c) \ge ov(c) \ \text{ for every branch node } c .
\]
\end{definition}

Every clean network is escape-free and the converse fails; the gap is not
cosmetic, since Batcher's 12-sorter \cite{Ba68} has a pass-through yet is
escape-free.

\section{The audit}
\label{sec:audit}

\subsection{The chain, isolated}

The section of \cite{VV72} bounding $P(2,N)$ opens with a premise, attributed to
the \emph{Transactions} note \cite{VV72a}:

\begin{quote}
``It has been observed \dots\ that the MAX subnetwork of any $N$-sorter is a
binary tree with $N$ leaves (the input leads) and $N-1$ branch nodes (the
comparators) rooted at $o_N$.'' \hfill --- p.~121

``Let $T$ be any $N$-sorter. The MAX subnetwork of $T$, $\MAX(T)$, includes
$N-1$ comparators, which we label $c_1, c_2, \dots, c_{N-1}$.''
\hfill --- p.~121
\end{quote}

\noindent
For a branch node $c_j$, $\nc(c_j)$ counts the comparators from the two inputs
of its subtrees to $c_j$ and on to $o_N$, comparator $c_j$ itself counted once;
$q_j$ is the second maximum's onward path to $o_{N-1}$, and the chapter asserts
that the paths $q_j$ ``together form a binary tree rooted at $o_{N-1}$, which we
call the MAX2 subnetwork'' (p.~122). Then

\begin{quote}
``\dots\ the two paths \dots\ together include $\nc(C_j) + \nc(q_j)$
comparators. Therefore, $p(2,T) = \max_j\,[\nc(C_j) + \nc(q_j)]$.''
\hfill --- p.~122, eq.~(5)

``$\MAXTWO(T)$ is a binary tree, so the path lengths $\nc(q_j)$ satisfy
$\sum_j 2^{-\nc(q_j)} = 1$.'' \hfill --- p.~124, eq.~(6)
\end{quote}

\noindent
From (5) and (6) alone the chapter derives (7),
$\sum_j 2^{-[p(2,T)-\nc(C_j)]} \le 1$, hence (8),
$p(2,T) \ge \lceil \log_2 f(\MAX(T))\rceil$, hence
$P(2,N) \ge \lceil \log_2 F(N)\rceil$ and the tabulated value $35+9$ at $N=13$.
There are no side conditions anywhere in the section: no standard-form
restriction, no normalization, no minimality hypothesis, no restriction on $N$.
Step (7) needs only the inequalities $\nc(q_j) \le p(2,T) - \nc(C_j)$ and
$\sum_j 2^{-\nc(q_j)} \le 1$; the failures below are in the fatal direction
anyway.

\subsection{The counterexample}

\begin{proposition}[the 3-sorter refutes (5) and (6)]
\label{prop:cex}
Let $T_1 = [(0,1),(0,2),(1,2)]$, the Batcher 3-sorter, with $|T_1| = 3 = S(3)$.
Then $|\MAX(T_1)| = 3$, not $N-1 = 2$; equation (5) returns $4$ for a quantity
whose value is $3$; equation (6)'s sum is $5/4$; and under the literal reading
of the chapter's equation (9), equation (8) itself is false on $T_1$.
\end{proposition}

\begin{proof}
Write $c_0, c_1, c_2$ for the three comparators in order. Every one of them is
traversed by a max-path, so $\MAX(T_1) = \{c_0,c_1,c_2\}$; but
$c_1 = (0,2)$ is a pass-through, since the maximum enters and leaves it on
channel~2 and no max-path arrives on channel~0 there. The branch nodes are
$\{c_0, c_2\}$ --- two of them, as the chapter says. What fails is their
identification with the comparators of $\MAX(T)$.

Take $c_0$, with the input pair (channel 0, channel 1). The maximum's path is
$c_0, c_2$ and the second maximum's is $c_0, c_1, c_2$, so
$\nc(c_0) = \nc(q_0) = 2$ and (5) asserts that the two paths together include
$4$ comparators. Their union is $\{c_0,c_1,c_2\}$: \emph{three} --- after
separating at $c_0$ the two values \emph{meet again} at $c_2$. Since
$|T_1| = 3$, we have $p(2,T_1) = 3$, confirmed by brute force. So (5)
\emph{exceeds} $p(2,T)$, which is exactly the direction that invalidates (7).

The claimed MAX2 leaves are the low output leads of $c_0$ and $c_2$, at depths
$2$ and $0$ --- the low output of $c_2$ \emph{is} $o_{N-1}$ --- giving
$\sum_j 2^{-\nc(q_j)} = 5/4$. The cause is structural: $q_0$ passes through
$c_2$'s low output lead, so the claimed leaves are not an antichain and
$\MAXTWO(T_1)$ is not a tree. Finally, with $\nc$ read literally as defined on
p.~122, $f = 2^2 + 2^3 = 12$ and (8) asserts $p(2,T_1) \ge 4 > 3$; with $f$
computed by the chapter's Theorem~1 on the branch tree, $f = 8$ and (8) holds
with equality.
\end{proof}

Both audit verifiers recompute this independently, and it is checkable by hand.
The \emph{tree} reading of $\nc$ is forced by the chapter's Fig.~5, whose
printed values are exactly the branch-tree quantities, and it is the reading
under which the chapter's conclusion is empirically unrefuted; it rescues
neither (5) nor (6). The exact replacement for (5) is
$|W(x,y)| = \nc(c) + \nc(q_c) - ov(c)$, and equation (5) is the claim that
$ov(c) = 0$.

\subsection{The failure is common, and not confined to slack}

Over $387$ constructed sorters (bubble, insertion, odd-even transposition and
Batcher for $n = 3,\dots,10$; thinned random prefixes; random channel
relabellings), each verified to be a sorter by the 0/1 principle
(\texttt{verify\_kraft\_dispute.py}, seed \texttt{20260818}):

\begin{center}
\begin{tabular}{lrr}
\toprule
statement & count & share \\
\midrule
$\MAX(T)$ has a pass-through & 236 & $61.0\,\%$ \\
some $q_j$ re-meets the max-path after $c_j$ & 233 & $60.2\,\%$ \\
eq.~(5) right-hand side strictly exceeds $p(2,T)$ (fatal) & 172 & $44.4\,\%$ \\
eq.~(6) sum $>1$ (fatal) & 164 & $42.4\,\%$ \\
eq.~(8) fails, literal $\nc$ & 102 & $26.4\,\%$ \\
eq.~(8) fails, \textbf{tree} $\nc$ & \textbf{0} & $0.0\,\%$ \\
conclusion $p(2,T) \ge \lceil \log_2 F(n)\rceil$ fails & \textbf{0} & $0.0\,\%$ \\
tight conclusion \emph{and} eq.~(6) broken & \textbf{105} & $\mathbf{27.1\,\%}$ \\
\bottomrule
\end{tabular}
\end{center}

The last row answers the natural dismissal, that this is a technicality biting
only where the bound has slack. The exhaustive check is sharper. Over the
complete set of $149{,}040$ size-optimal 5-sorters, equation (8) is tight on
every single one, and equation (6)'s sum exceeds 1 on $43.9\,\%$ of them
(\texttt{verify\_kraft\_wave1.py}, \texttt{verify\_kraft\_wave2.py}): the broken
step fails on nearly half of a population in which the conclusion has no slack
at all. Nor is the violation bounded --- deterministic adversarial search
reaches Kraft sums of $1.6875$ at $n=6$ and $1.8125$ at $n=7$, so the suprema
rise with $n$, and no argument may assume a constant bound on the excess. The
percentages over constructed populations are properties of the generators and
seeds, not rates (Section~\ref{sec:trust}); the shipped scripts \emph{assert}
the existence and universality statements and \emph{print} the counts.

\subsection{The strongest defence, and its refutation}

Ten steelman defences were attempted under an adversarial brief and are recorded
with their refutations in \cite{Art}. Two matter here. The first is the
slack argument just refuted. The second is the one a sympathetic reader
reconstructs first, and, we believe, what the chapter had in mind:

\begin{quote}
\textbf{$\MAX(T)$ and $\MAXTWO(T)$ are comparator-disjoint.} Every comparator of
$\MAX(T)$ has both inputs on MAX-edges; $q_j$ starts on $c_j$'s \emph{lower}
output, which is not a MAX-edge; a comparator with a non-MAX input is not in
$\MAX(T)$; induct.
\end{quote}

Disjointness would make (5) an exact union with nothing double-counted, and
would force every $q_j$ onto high leads throughout --- hence prefix-free, hence
Kraft --- repairing (5) and (6) at one stroke. It fails because its premise is
precisely what a pass-through violates: a pass-through lies in $\MAX(T)$ and has
one blue input. On $T_1$ the two subnetworks share two comparators, one of them
the root branch node; across the sweep they meet in $60.2\,\%$ of networks and
at a branch node in $59.9\,\%$. Disjointness \emph{does} hold on Batcher's
8-sorter, which is why the chapter's own worked example cannot expose it. The
conclusion itself survives every test we have: with no pruning at all, all $42$
sorting sequences on 3 channels of length $\le 4$ and all $912$ on 4 channels of
length $\le 6$ satisfy it (\texttt{verify\_kraft\_dispute.py}).

\subsection{A misstatement is not a gap}

The chapter's only worked example is Batcher's odd-even 8-sorter. Reconstructed
from the algorithm rather than from any table, its $\MAX(T)$ has seven
comparators, seven branch nodes and \emph{no} pass-through; its $\nc$ multiset
$\{3,3,3,3,4,4,5\}$ matches the values printed in the chapter's Fig.~5, and
$f = 96 = F(8)$ (\texttt{verify\_huffman2.py}). But its MAX2 depths give
$\sum_j 2^{-\nc(q_j)} = 0.75 \ne 1$, so equation (6) is not even an equality in
the chapter's own example. The two findings must be kept apart, since conflating
them would overstate the case: $0.75$ is \emph{below} 1, step (7) needs only
$\le 1$, and the derivation goes through there unharmed, whereas $T_1$ gives
$1.25$. This also suggests how the error survived --- the single example in the
chapter fails the equation on the safe side, and its MAX subnetwork is clean, so
it discriminates neither the literal from the tree reading of $\nc$ nor a
correct premise from a false one.

\section{The record: twenty entries in three inconsistent tables}
\label{sec:record}

\subsection{No erratum, and no reader}

We searched for any erratum, correction, gap note, reproof or formalization of
equations (5)--(9): the OpenAlex and Google Scholar citation graphs, arXiv title
and full-text search, DBLP, GitHub code and issue search, Knuth's Vol.~3 and its
errata, the Stanford technical-report indexes including the CS-TR scan
collection released in December 2025, Dobbelaere's table with its Wayback
history, and Harder's Isabelle sources. Nothing exists, and the reason is
stronger than ``nobody found the error'': nobody has engaged with the argument.
A search for \texttt{P(2,}, \texttt{MAX2}, \texttt{MAX subnetwork} and
\texttt{Kraft} across some twenty-five modern papers returns zero hits, and
Knuth carries no two-value exercise and never cites the chapter. There is no
machine-checked version of the $P(2,N)$ bound.

\subsection{Twenty entries, not one}

The chapter's chain is not used at $N=13$ only. Taking
\[
  S(N) \;\ge\; \max\bigl(\,S(N-1) + \lceil \log_2 N\rceil,\;\;
  S(N-2) + \lceil \log_2 F(N)\rceil \,\bigr),
  \qquad S(11) = 35,
\]
and regenerating $F(N)$ to $N=32$ by exact minimization over shapes --- not by
the chapter's relaxation (19), which is known to agree only for $N \le 16$ ---
reproduces the published column exactly at every one of $N = 13,\dots,32$. The
regeneration is cross-checked two ways: it reproduces the chapter's Table~1 for
$N \le 16$, and at $N = 13$ brute-force enumeration of all $208{,}012$ shapes
with 13 leaves agrees with the height dynamic program
(\texttt{verify\_novel\_facts.py}).

\begin{center}
\footnotesize
\begin{tabular}{r r r c c c c r}
\toprule
$N$ & $F(N)$ & $\lceil\log_2 F\rceil$ & one-value (proved) & eq.~(8) chain
& Dobbelaere \cite{Do25} & Wikipedia \cite{Wik26} & deficit \\
\midrule
13 & 392 & 9 & \textbf{43} & 44 & 44 & 43 & 1 \\
14 & 424 & 9 & \textbf{47} & 48 & 48 & 47 & 1 \\
15 & 480 & 9 & \textbf{51} & 53 & 53 & 51 & 2 \\
16 & 512 & 9 & \textbf{55} & 57 & 57 & 55 & 2 \\
17 & 784 & 10 & \textbf{60} & 63 & 63 & 60 & 3 \\
18 & 848 & 10 & \textbf{65} & 68 & 68 & 65 & 3 \\
19 & 960 & 10 & \textbf{70} & 73 & 73 & 70 & 3 \\
20 & 1024 & 10 & \textbf{75} & 78 & 78 & 75 & 3 \\
21 & 1232 & 11 & \textbf{80} & 84 & 84 & --- & 4 \\
22 & 1296 & 11 & \textbf{85} & 89 & 89 & --- & 4 \\
23 & 1408 & 11 & \textbf{90} & 95 & 95 & --- & 5 \\
24 & 1472 & 11 & \textbf{95} & 100 & 100 & --- & 5 \\
25 & 1872 & 11 & \textbf{100} & 106 & 106 & --- & 6 \\
26 & 1936 & 11 & \textbf{105} & 111 & 111 & --- & 6 \\
27 & 2048 & 11 & \textbf{110} & 117 & 117 & --- & 7 \\
28 & 2112 & 12 & \textbf{115} & 123 & 123 & --- & 8 \\
29 & 2320 & 12 & \textbf{120} & 129 & 129 & --- & 9 \\
30 & 2384 & 12 & \textbf{125} & 135 & 135 & --- & 10 \\
31 & 2496 & 12 & \textbf{130} & 141 & 141 & --- & 11 \\
32 & 2560 & 12 & \textbf{135} & 147 & 147 & --- & 12 \\
\bottomrule
\end{tabular}
\end{center}

Three things are in that table. First, it is the first written derivation of
Dobbelaere's 2025 lower-bound column. That page carries the numbers with a
one-line changelog --- ``2025-04-21 Tighter lower bounds for size, on suggestion
of Jelmer Firet and based on principles in [VVoorh72]'' --- no derivation, no
citation on the row (the only reference attached to the $N=13$ row is for the
depth result), and we located no publication, preprint or repository by the
person credited. We can now say exactly what the column is, and we confirm it
constructively rather than by absence of evidence. Second, every one of the
twenty entries strictly exceeds the one-value floor, and the only route to the
excess anywhere in the literature is the equation whose published derivation is
invalid: the audit bears on $S(13)$ through $S(32)$. Third, the deficit column
is the cost of the error, and it grows --- one comparator at $N=13$, twelve at
$N=32$.

\subsection{Which column is defensible}

The two most consulted public tables disagree at every $N$ from 13 to 20.
Wikipedia's ``Size, lower bound'' row is exactly the one-value column, and
footnotes itself ``Obtained by Van Voorhis lemma and the value $S(11) = 35$''
\cite{Wik26}; Dobbelaere's is the eq.~(8) column \cite{Do25}; and OEIS A003075
\cite{OEIS}, which Dobbelaere's column header credits, lists no lower bound for
any $N \ge 13$ and has no b-file. So the record is three-way inconsistent, and
this work adjudicates it: the defensible column is Wikipedia's.

\begin{quote}
The published lower-bound table for sorting-network size has no unconditional
support at any $N \ge 13$. The unconditional table for $N = 13,\dots,32$ is\\
$43,\, 47,\, 51,\, 55,\, 60,\, 65,\, 70,\, 75,\, 80,\, 85,\, 90,\, 95,\,
100,\, 105,\, 110,\, 115,\, 120,\, 125,\, 130,\, 135$.
\end{quote}

\noindent
The twenty numbers under audit are not refuted, and we make no claim that any is
false; the claim is exactly and only about the status of their proof. Nor is a
peer-reviewed result overturned --- the value at $N = 13$ has never appeared in
a peer-reviewed publication, and the best published lower bound is 43 --- which
lowers the stakes and correspondingly raises the value of recording the finding,
since the column is propagating with no correct argument behind it.

\subsection{What the table rests on}

Both columns reduce to a single computed value and to no other: $S(12) = 39$ is
one-value from $S(11) = 35$, and every entry at $N \ge 13$ is a chain above
$S(11)$ or $S(12)$; replacing $S(11) = 35$ by 34 lowers every entry at
$N \ge 12$ by at least one (\texttt{verify\_novel\_facts.py}). That value had
exactly one machine proof, Harder's 2.9\,GB certificate \cite{HaCert}. A
companion paper reports a second, independent one, from a rebuilt search engine
whose certificate is accepted by an unmodified extraction of the same
Isabelle/HOL-verified checker \cite{P2,Cert}: the deepest computational input to
every published lower bound above $n = 11$ has now been confirmed twice.

\section{The partial repair}
\label{sec:repair}

The gap is the missing antichain step. We supply it for two classes, give a
complementary bound at the opposite extreme, close one family of repairs
structurally, and then say how much that leaves.

\subsection{The exact correction of equation (5)}

\begin{lemma}[exact decomposition]
\label{lem:decomp}
Let $c = \LCA_B(x,y)$ and let $e(z,c)$ be the number of comparators strictly
between input $z$ and $c$ on $z$'s max-path. Then
$|W(x,y)| = e(x,c) + e(y,c) + g(c) + r(c)$, where $g(c)$ is the number of
comparators from $c$ to $o_N$ inclusive on the maximum's onward path and
$r(c) = nq(c) - ov(c)$ is the number the second maximum traverses after $c$
that the maximum does not. In particular $g$ and $r$ depend only on $c$.
\end{lemma}

The proof is three short observations, all general and all in \cite{Art}: the
two values first meet at $\LCA_B(x,y)$; after $c$ their trajectories are
functions of the pair of occupied leads alone; and the two pre-meeting paths are
disjoint. This is the correction of equation (5), which double-counts $ov(c)$.
Aggregating over the admissible pairs at a node gives the identity that replaces
it.

\begin{lemma}[the identity]
\label{lem:identity}
For every $N$-sorter and every branch node $c$,
\[
  W^*(c) \;=\; a(c) + \varepsilon(c) + \gamma(c) + r(c),
  \qquad \varepsilon, \gamma, r \ge 0 ,
\]
where $\varepsilon(c)$ measures how far the maximum's pre-meeting paths exceed
the branch-tree heights and $\gamma(c)$ is exactly the number of pass-throughs
on the maximum's stem from $c$ to $o_N$. Hence
$p(2,T) = \max_c [\,a(c) + \varepsilon(c) + \gamma(c) + r(c)\,]$.
\end{lemma}

This is an exact expression for $p(2,T)$ in branch-tree data plus named
non-negative corrections, verified with zero violations over the
$32{,}043$-sorter corpus of \texttt{verify\_kraft\_wave2.py}. One consequence
deserves separate statement: equation (8) is \emph{equivalent} to
\[
  \sum_c 2^{a(c)} \;\le\; 2^{\,\max_c [\,a(c)+\varepsilon(c)+\gamma(c)+r(c)\,]},
\]
a statement about a binary tree decorated with three non-negative integer
functions, and it is \emph{false} for arbitrary decorations, since
$\varepsilon = \gamma = r = 0$ refutes it on any tree with two or more internal
nodes. The entire content of equation (8) is therefore a realizability
constraint --- which decorated trees arise from an actual sorter --- and that is
the shape of the missing lemma.

\subsection{The escape-free case}

\begin{theorem}[equation (8) for escape-free sorters]
\label{thm:ef}
For every escape-free $N$-sorter $T$ with branch tree $B$,
\[
  1 \;\ge\; \sum_c 2^{-nq(c)} \;\ge\; \sum_c 2^{-(p(2,T)-a(c))}
    \;=\; 2^{-p(2,T)} f(B) ,
\]
hence $p(2,T) \ge \lceil \log_2 f(B)\rceil$ and
$|T| \ge S(N-2) + \lceil \log_2 f(B)\rceil$.
\end{theorem}

The proof is elementary and is in \cite{Art}, its lemmas checked by
\texttt{verify\_kraft\_wave1.py} and \texttt{verify\_kraft\_wave2.py}. The two
steps the chapter omits are these. In the colouring of Definition~\ref{def:one},
every low output lead is blue, a high output lead is red exactly when its
comparator has a red input, and a comparator is a branch node exactly when both
its inputs are red. In an escape-free network the second maximum therefore
traverses \emph{no branch node} after the meeting, so the $N-1$ sources --- the
low output leads of the branch nodes --- are pairwise non-ancestral under the
second maximum's lead successor map. That is the antichain, and Kraft's
inequality applies; the per-node inequality $nq(c) \le p(2,T) - a(c)$ is
Lemma~\ref{lem:decomp} with the deepest-leaf pairing. Two things change relative
to the chapter, both essential: equation (6) becomes an \emph{inequality}, which
is why Batcher's own 8-sorter gives $0.75$ and the argument is unharmed, and the
antichain is \emph{proved} rather than assumed. In an arbitrary sorter the
lead-level successor map is not even a function, which is the precise sense in
which $\MAXTWO(T)$ is not a tree.

\subsection{Isolating the hypotheses}

Escape-freeness enters Theorem~\ref{thm:ef} only through (C) and (D) of
Definition~\ref{def:two}. Isolating them widens the class.

\begin{theorem}[equation (8) under (C) and (D)]
\label{thm:cd}
If an $N$-sorter satisfies (C) and (D), then
$p(2,T) \ge \lceil \log_2 f(B)\rceil$.
\end{theorem}

\emph{Proof:} verbatim that of Theorem~\ref{thm:ef}, which uses escape-freeness
nowhere else. The gain is small and we say so. On the $32{,}043$-sorter corpus
(C)$\wedge$(D) holds on $7{,}214$ networks against $7{,}167$ escape-free ones, a
widening of $0.7\,\%$; over the exhaustive $n = 5$ universe it is $33{,}030$
against $32{,}850$, a widening of $0.55\,\%$; and at $n = 4$, in both universes,
(C) and (D) are disjoint on escaping sorters, so the conjunction buys nothing at
all. Equation (8) holds on all $7{,}214$. The value of Theorem~\ref{thm:cd} is
not its coverage but that it names the two things, and only the two things, a
general proof must supply.

\subsection{The pass-through-rich end}

\begin{proposition}[an upper bound on $f(B)$]
\label{prop:rich}
For every $N$-sorter, $f(B) \le 2^{p(2,T)} \sum_c 2^{-\gamma(c)}$. In
particular equation (8) holds whenever $\sum_c 2^{-\gamma(c)} \le 1$, and
whenever the root's stem carries at least one pass-through and every edge of $B$
carries at least two.
\end{proposition}

Unlike (C) and (D), this hypothesis is computable from the MAX subnetwork alone,
and the class is non-empty and not contained in the escape-free class.
\textbf{But it must be quoted with its caveat.} Appending inert comparators
raises $p(2,T)$ while leaving $f(B)$ fixed, so the sorters it reaches are
exactly those with large slack --- the witnesses have $p(2,T) = 7$ against a
bound of $5$. It is recorded because it completes the picture from the opposite
direction, not because it is useful for size-optimal networks.

\subsection{Coverage: where the open problem lives}

This is the most useful new fact here, and the most uncomfortable. Over the
$32{,}043$-sorter corpus of \texttt{verify\_kraft\_wave2.py} --- exhaustive at
$n = 4$ in both universes and at $n=5$ in the minimal universe, plus constructed
families to $n = 11$, inert-augmented families and random $n = 6,7$ sorters ---
each condition fires as follows.

\begin{center}
\begin{tabular}{lrrl}
\toprule
sufficient condition & count & fraction & kind \\
\midrule
clean & 5{,}891 & 0.1838 & structural \\
escape-free (Theorem~\ref{thm:ef}) & 7{,}167 & 0.2237 & structural \\
hypothesis of Proposition~\ref{prop:rich} & 0 & 0.0000 & structural \\
\textbf{structural union} & \textbf{7{,}167} & \textbf{0.2237} & \\
(C)$\wedge$(D) (Theorem~\ref{thm:cd}) & 7{,}214 & 0.2251 & computed \\
$\sum_c 2^{-\gamma(c)} \le 1$ & 18 & 0.0006 & computed \\
$\sum_c 2^{-r(c)} \le 1$ & 16{,}052 & 0.5010 & computed \\
\textbf{union of all six} & \textbf{16{,}053} & \textbf{0.5010} & \\
\bottomrule
\end{tabular}
\end{center}

Equation (8) fails on none of the $32{,}043$. Structural and computed conditions
must not be blurred: escape-freeness and the hypothesis of
Proposition~\ref{prop:rich} are properties of a network's shape, whereas (C),
(D) and the two sums are numbers one computes per network, and a per-network
computation is not a theorem. The proved classes cover $22.4\,\%$ of the corpus
structurally, and \textbf{$49.9\,\%$ of it has no proof of equation (8) by any
route in this work}. Proposition~\ref{prop:rich}'s hypothesis fires on no
network here at all: its witnesses have to be built by inert padding, which is
exactly the slack regime its caveat describes. The uncovered zone is
\emph{intermediate} in pass-through count. On this corpus the fraction with no
covering condition is $0.000$ at zero pass-throughs --- all $5{,}891$ such
networks are clean, hence covered --- and then $0.346$, $0.740$ and $0.995$ at
one, two and three; beyond three the corpus thins to a few dozen networks per
bucket and the fraction falls back irregularly, reaching zero only at the
largest counts present. Equation (8) is provable by Theorem~\ref{thm:ef} at the
pass-through-poor end and by Proposition~\ref{prop:rich} at the
pass-through-rich end, and the open problem lives strictly in between.

\subsection{What is closed off}

\begin{proposition}[per-node charging is impossible]
\label{prop:pernode}
On any $T$ with $f(B) = 2^{p(2,T)}$ one has $\sum_c 2^{a(c)-p(2,T)} = 1$
exactly. Hence for any per-node exponent $X(c) \le p(2,T)$,
$\sum_c 2^{a(c)-X(c)} \ge 1$, with equality if and only if $X(c) = p(2,T)$ at
every branch node --- and tight sorters with non-constant $W^*(c)$ exist. No
certificate built from each node's own pruning can prove equation (8).
\end{proposition}

The witness is an 18-comparator 7-sorter with $f = 2^{p(2,T)} = 128$ and
$W^* = \{6,7,7,7,7,7\}$; the strongest member of the family, charging each node
against its own $W^*(c)$, is refuted outright by an adversarially found
15-comparator 6-sorter with $\sum_c 2^{a(c)-W^*(c)} = 33/32$. Four further
repair families are closed with explicit counterexamples in \cite{Art}. What
remains is worth stating plainly: without the Kraft step the chain yields only
$p(2,T) \ge \max_j \nc(c_j)$, which for the $f$-minimizing 13-leaf trees is
$42$. The Kraft step carries the entire excess over the one-value floor, and the
Kraft step is the broken one.

\subsection{One thing the repair does buy}

Theorem~\ref{thm:ef} is per-network and un-minimized, over the network's own
branch tree --- exactly the hypothesis a shape-based case analysis of the
$n = 13$ problem has to assume outright. Two machine checks close the
identification: $f(B)$ equals the shape-theoretic $V(T)$ on all $208{,}012$
plane shapes with 13 leaves and on every shape with at most 11 leaves, and at
$N = 13$ and the tabulated size Theorem~\ref{thm:ef} forces $f(B) \le 512$,
which is that analysis's admissibility condition
(\texttt{verify\_novel\_facts.py}, \texttt{verify\_huffman2.py}).

\begin{corollary}[the shape case split, unconditional on the escape-free class]
\label{cor:shapes}
Let $C$ be an escape-free 13-sorter of the tabulated size. Then its max branch
tree is one of exactly $84$ admissible plane shapes ($6$ up to reflection), its
root splits the 13 leaves as $5+8$, $6+7$ or $4+9$, and its leaf-depth profile
is one of exactly four. If $C$ and its reverse-and-flip dual are both
escape-free, both branch trees are so constrained.
\end{corollary}

The finite shape enumeration and the identification $f = V$ are machine-checked;
Corollary~\ref{cor:shapes} inherits from Theorem~\ref{thm:ef} the status of a
human proof. What changes is that the case split no longer needs the chapter: it
is unconditional on a class we can define, check and measure.

\section{What is not disturbed}
\label{sec:undisturbed}

\subsection{The one-value bound, and Harder's Lemma 17}

A reader who checks Harder's paper will find the same identification we attack.
His Lemma 17 reproduces the classical one-value bound with a full proof, and in
it:

\begin{quote}
``if we take the union of the pruned paths for all $i$ and remove the common
output $n$, we obtain a binary tree rooted in the comparator gate connected to
output $n$, where the leaves are all inputs of $c$ and all inner vertices are
comparator gates \dots\ the leaf for every input $i$ has a depth of
$\delta(c,i)$.'' \hfill --- \cite{Ha20}, proof of Lemma 17
\end{quote}

Identifying max-path comparators with branch nodes is exactly the p.~121
premise, and by Proposition~\ref{prop:cex} it is false whenever pass-throughs
exist: $\delta(c,i)$ is a literal comparator count and can exceed the
branch-tree depth. \textbf{It is harmless there}, for a one-line reason. The
one-value bound needs only
$\max_i \delta(c,i) \ge \max_i \mathrm{depth}_B(i) \ge \lceil\log_2 n\rceil$,
and literal depth \emph{dominates} branch depth, so the inequality goes the safe
way. $S(11) = 35$, $S(12) = 39$ and the published $S(13) \ge 43$ stand. We
record this because a reader who notices Harder's phrasing will otherwise ask
whether the audit also breaks $S(12)$, and because it strengthens the account of
how the error survived: the same imprecision was reproduced, unnoticed, in a
2020 paper that is the current authority for $S(11)$ and $S(12)$. One smaller
point. Harder's Lemma 17 is a paper proof --- his Isabelle development verifies
the certificate checker, not Van Voorhis's bound --- so $S(12) = 39$ is not
machine-checked either.

\subsection{What is still known about a minimum-size 13-sorter}

The following is unconditional. It uses only that deleting a channel at either
polarity from an $N$-sorter removes exactly the comparators on that channel's
extremal path and leaves an $(N-1)$-sorter, so $|C| \ge S(N-1) + \delta(C,i)$;
that $\delta(C,i) = \mathrm{depth}_B(i) + pt(C,i)$, which is
Proposition~\ref{prop:cex} used in the safe direction; Kraft equality on a full
binary tree; and the duals for the minimum. No step uses equation (8). Write $m$
for the tabulated size.

\begin{proposition}[unconditional structure]
\label{prop:struct}
Let $C$ be a 13-sorter with $|C| = m$. Then:
(i) every extremal path of $C$, at either polarity, has at most $m - S(12) = 5$
comparators, so $\mathrm{depth}_B(i) + pt(C,i) \le 5$ for every channel $i$;
(ii) a leaf at branch-depth 5 has a pass-through-free max-path;
(iii) both branch trees have height in $\{4,5\}$, and the leaf-depth multiset of
either is one of exactly 13 Kraft-tight profiles with all depths at most 5;
(iv) if any extremal path has length 5, deleting that channel leaves a
size-optimal 12-sorter in which both branch trees have height exactly 4;
(v) $C$ has a one-channel deletion minor that is a 12-sorter with 39 or 40
comparators and a two-channel deletion minor that is an 11-sorter with 35 or 36,
and these are the only two places in the deletion chain where the slack is 1;
and (vi) $|\mathrm{maxpath}(i) \cup \mathrm{minpath}(j)| \le m - S(11) = 9$ for
all $i \ne j$.
\end{proposition}

Item (ii) is the point of contact: it links pass-throughs, the audit's central
object, to the $n = 13$ question directly and without appeal to the disputed
inequality. Item (vi) is the max/min analogue of the chapter's max/second-max
quantity, and it is worth a sentence because of what it avoids. The chapter's
two-channel deletion removes the maximum and the \emph{second} maximum, and its
Kraft step needs $\MAXTWO(T)$ to be a binary tree, which is the false premise;
the max/min deletion removes the maximum and the \emph{minimum}, and both
objects it involves are honest Kraft-tight binary trees with $N$ leaves. Whether
$\min_C \max_{i \ne j} |\mathrm{maxpath}_C(i) \cup \mathrm{minpath}_C(j)|$ admits
a Kraft-style lower bound we do not know --- the two paths can share comparators
--- nor have we established that the variant is new.

\section{What we trust}
\label{sec:trust}

Following Harder \cite{Ha20}, we enumerate the trusted base rather than claim
correctness. We do not require, or expect, it to be free of bugs; only that no
bug was triggered in the runs reported here.

\begin{center}
\small
\begin{tabular}{p{0.29\textwidth} p{0.14\textwidth} p{0.44\textwidth}}
\toprule
component & trusted base? & note \\
\midrule
hardware, operating system, Python 3.14 & \textbf{yes} & general-purpose \\
the two audit verifiers of Section~\ref{sec:audit} & \textbf{yes} & mutually
independent, no shared code; the counterexample needs neither and is checkable
by hand \\
the two wave verifiers of Section~\ref{sec:repair} & \textbf{yes} & they import
the audit verifier's primitives, so they are extensions, not independent
re-derivations \\
Theorems \ref{thm:ef}, \ref{thm:cd}; Lemmas \ref{lem:decomp},
\ref{lem:identity}; Propositions \ref{prop:rich}, \ref{prop:pernode},
\ref{prop:struct} & \textbf{yes} & human proofs; no proof assistant has checked
them \\
the identification $f = V$ and the 13-leaf shape enumeration & \textbf{yes} &
our own code, run exhaustively \\
the prevalence and coverage percentages & \textbf{yes} --- scoped &
properties of the named corpora and seeds, not rates \\
Isabelle/HOL kernel and extraction; GHC & \textbf{yes} & for $S(11) = 35$
only \\
the search engines and certificate generators behind $S(11) = 35$ & no & the
value is re-derived by the verified checker \cite{HaCert,Cert,P2} \\
the two archival items \cite{VV71,VV72a} & \textbf{yes} --- not obtained &
see below \\
\bottomrule
\end{tabular}
\end{center}

Five things follow, and none of them is a formality.

\textbf{Independence is not uniform.} The two audit scripts share no code, so
the counterexample and the prevalence figures are re-derived twice from
scratch. But the two wave verifiers that carry Section~\ref{sec:repair}
\emph{import} the first script's primitives; they are extensions, not
independent re-derivations, and should be read as such.

\textbf{We do not meet the de Bruijn criterion for the new mathematics.} A
computer-assisted proof should be checkable by an independent small program.
The certified value $S(11) = 35$ meets that criterion. Theorems
\ref{thm:ef} and \ref{thm:cd} do not: they are human proofs whose lemmas are
tested on samples, and Proposition~\ref{prop:pernode}'s own history --- 900
samples with supremum exactly $1.000000$, followed immediately by an adversarial
counterexample --- is why we do not offer sample agreement as more than weak
evidence. Formalizing Theorem~\ref{thm:ef} is a small task and is the natural
next step; it would also make Corollary~\ref{cor:shapes} fully machine-checked.

\textbf{Every percentage here is a corpus property, not a rate.} The
$61.0\,\%$, $44.4\,\%$ and $27.1\,\%$ of Section~\ref{sec:audit} are properties
of a fixed generator set and seed, and the coverage fractions are properties of
the $32{,}043$-sorter corpus. Both are regenerated by the named scripts, and
both establish that the phenomena are common in ordinary networks rather than
that they occur at any rate under a measure on sorting networks. The two
exhaustive populations --- all $708$ minimal 4-sorters and all $149{,}040$
minimal 5-sorters --- are the only figures here that are rates, and they are
rates only at those two widths. We do not know what fraction of
\emph{size-optimal} networks at large $n$ is escape-free, and a reader is
entitled to ask.

\textbf{Two archival items were not obtained, and one may be the origin of the
error.} The \emph{Transactions} note \cite{VV72a} is what the chapter cites
specifically for the false structural claim. Three independent readers of it ---
Knuth, Parberry \cite{Pa89} and Harder --- report only the one-value bound, and
Harder's rendition contains the same identification, which makes it likely the
false premise originates there; if the note states the claim correctly, the
chapter's error is introduced rather than inherited. The 1971 Stanford
dissertation \cite{VV71} is not digitized and is absent from the Stanford CS-TR
scan collection released in December 2025, which we checked; only the
divide-sort-merge report \cite{VV71b} is there. Obtaining either is
procurement, not research, and neither affects $S(13) \ge 43$.

\textbf{The status of the disputed inequality.} Equation (8) is neither proved
nor refuted. It has substantial empirical support --- zero violations across the
constructed and exhaustive corpora above, roughly 15{,}000 adversarially
searched networks aimed directly at it \cite{Art}, and consistency with every exact $S(n)$
known, the slacks $S(N) - S(N-2) - \lceil\log_2 F(N)\rceil$ for $N = 3,\dots,12$
being $0,0,0,1,0,0,1,2,2,1$, so a single negative entry would have refuted it. A
full repair must be global, must \emph{handle} pass-throughs rather than exclude
them, must survive the tight regime, and must reach $\lceil\log_2 f(B)\rceil$
rather than $\max_j \nc(c_j)$. Failing that, the honest state of the record is
$43 \le S(13) \le 45$, and correspondingly at every $N$ up to 32. The audit was
carried out by language-model agents under a fixed protocol, one reading the
primary source and a second instructed to defend the chapter and refute the
first; no claim here rests on a model's assertion, since every quantitative
statement is produced by one of the five scripts and the central counterexample
is checkable on paper.

Van Voorhis closed his chapter by predicting that the next improvement in the
lower bound for $S(N)$ would come from an approximation of $P(3,N)$. Fifty-four
years on, $P(2,N)$ is still open.

\begin{thebibliography}{99}
\footnotesize
\setlength{\itemsep}{1pt plus 0.5pt}
\bibitem{Ba68} K. E. Batcher. Sorting networks and their applications.
\emph{AFIPS Spring Joint Computer Conference}, 1968, 307--314.

\bibitem{BN62} R. C. Bose, R. J. Nelson. A sorting problem.
\emph{J. ACM} 9(2), 1962, 282--296.

\bibitem{BZ14} D. Bundala, J. Z\'avodn\'y. Optimal sorting networks.
\emph{LATA 2014}, LNCS 8370, 236--247. arXiv:1310.6271.

\bibitem{BR26} R. Burca, M. Raschip. Finding minimal-size sorting networks using
deep Q-learning. \emph{ICAART 2026}.

\bibitem{CCFS16} M. Codish, L. Cruz-Filipe, M. Frank, P. Schneider-Kamp.
Sorting nine inputs requires twenty-five comparisons. \emph{J. Computer and
System Sciences} 82(3), 2016, 551--563. Conference version ICTAI 2014;
arXiv:1405.5754.

\bibitem{CCS14} M. Codish, L. Cruz-Filipe, P. Schneider-Kamp. The quest for
optimal sorting networks: efficient generation of two-layer prefixes.
\emph{SYNASC 2014}. arXiv:1404.0948.

\bibitem{CLS17} L. Cruz-Filipe, K. F. Larsen, P. Schneider-Kamp. Formally
proving size optimality of sorting networks. \emph{J. Automated Reasoning} 59,
2017, 425--454.

\bibitem{CS15} L. Cruz-Filipe, P. Schneider-Kamp. Formalizing size-optimal
sorting networks: extracting a certified proof checker. \emph{ITP 2015}.
arXiv:1502.05209.

\bibitem{CS24} L. Cruz-Filipe, P. Schneider-Kamp. Minimizing sorting networks at
the sub-comparator level. \emph{LPAR 2024}, EPiC Series in Computing.

\bibitem{Do25} B. Dobbelaere. \emph{Smallest and fastest sorting networks for a
given number of inputs.}
\url{https://bertdobbelaere.github.io/sorting_networks.html}; changelog entry of
2025-04-21; consulted 2026-08-22.

\bibitem{EM14} T. Ehlers, M. M\"uller. Faster sorting networks for 17, 19 and 20
inputs. arXiv:1410.2736, 2014.

\bibitem{FK73} R. W. Floyd, D. E. Knuth. The Bose--Nelson sorting problem. In
\emph{A Survey of Combinatorial Theory}, North-Holland, 1973, 163--172.

\bibitem{Ha20} J. Harder. \emph{An answer to the Bose--Nelson sorting problem
for eleven and twelve elements.} arXiv:2012.04400v3, 2022.

\bibitem{HaCert} J. Harder. \emph{Certificate for the minimal size of 11- and
12-channel sorting networks.} Zenodo, doi:10.5281/zenodo.4108365; software
\texttt{jix/sortnetopt}, doi:10.5281/zenodo.4139152.

\bibitem{Ju95} H. Juill\'e. Evolution of non-deterministic incremental
algorithms as a new approach for search in state spaces. \emph{ICGA 1995},
351--358.

\bibitem{KLM95} N. Kahale, T. Leighton, Y. Ma, C. G. Plaxton, T. Suel,
E. Szemer\'edi. Lower bounds for sorting networks. \emph{STOC 1995}, 344--353.

\bibitem{Kn73} D. E. Knuth. \emph{The Art of Computer Programming, Vol. 3:
Sorting and Searching.} Addison-Wesley, 1973. Section 5.3.4, exercise 42 and its
answer.

\bibitem{MG15} V. Marinov, D. Gregg. Sorting networks: the final countdown.
arXiv:1502.05983, 2015.

\bibitem{OEIS} OEIS Foundation Inc. Sequence A003075: minimal number of
comparisons needed for a sorting network on $n$ elements.
\emph{The On-Line Encyclopedia of Integer Sequences}, consulted 2026-08-22.

\bibitem{Pa89} I. Parberry. A computer-assisted optimal depth lower bound for
nine-input sorting networks. \emph{Supercomputing 1989}, 152--161.

\bibitem{Art} The authors. \emph{Artifact repository: verifiers, patches, theory
documents, evidence manifests.} To be deposited. Cited throughout for
check-level indexes, full proofs, refuted routes and measurement records.

\bibitem{Cert} The authors. \emph{Certificate of the minimal size of 11-channel
sorting networks, independently regenerated.} 2{,}442{,}317{,}348 bytes,
SHA-256 \texttt{672c433f\dots}. To be deposited.

\bibitem{P2} The authors. \emph{A cutoff theorem for exhaustive
sorting-network search, and an independently certified $S(11) = 35$.}
Companion paper.

\bibitem{TS26} A. Tozawa, K. Sadakane. Odd-even transposition sort is an optimal
stable standard sorting network. \emph{Information Processing Letters} 195,
2026, article 106659.

\bibitem{VV71} D. C. Van Voorhis. \emph{Efficient sorting networks.}
Ph.D. dissertation, Stanford University, 1971. Not obtained.

\bibitem{VV71b} D. C. Van Voorhis. \emph{A lower bound for sorting networks that
use the divide-sort-merge strategy.} Stanford Digital Systems Laboratory
TR 17, STAN-CS-71-238, 1971.

\bibitem{VV72} D. C. Van Voorhis. Toward a lower bound for sorting networks. In
R. E. Miller, J. W. Thatcher (eds.), \emph{Complexity of Computer Computations},
Plenum Press, 1972, 119--129. Equation numbers, Theorem 1, Table 1 and
Figs.~1--5 cited here are the chapter's own; page numbers are book pages.

\bibitem{VV72a} D. C. Van Voorhis. An improved lower bound for sorting networks.
\emph{IEEE Transactions on Computers} C-21(6), 1972, 612--613. Not obtained.

\bibitem{Vlad25} A. Vlad, B. Pascu, C. Cioata, M. Raschip. Identifying
optimal-size sorting networks with reinforcement learning. \emph{ICTAI 2025}.

\bibitem{Wik26} \emph{Sorting network.} Wikipedia, revision of 2026-08-15;
consulted 2026-08-22.

\end{thebibliography}

\end{document}
